\subsection{拓扑群的积}

设 \((\mathbf{G}_i)_{i\in I}\) 是一个拓扑群族。则我们可以在积集上定义一个群结构

\[
\mathbf {G} = \prod_ {i \in I} \mathbf {G} _ {i}
\]

（诸 \(\mathbf{G}_i\) 的群结构的积），办法是定义 \((x_i) \cdot (y_i) = (x_i y_i)\)。若 \(e_i\) 是 \(\mathbf{G}_i\) 的单位元，则 \(e = (e_i)\) 是 \(\mathbf{G}\) 的单位元，且有 \((x_i)^{-1} = (x_i^{-1})\)。诸 \(\mathbf{G}_i\) 的拓扑的积拓扑（第一章 § 2 第 3 目）与这个群结构相容。因为映射 \(((x_i), (y_i)) \to (x_i y_i^{-1})\)（\(\mathbf{G} \times \mathbf{G}\) 到 \(\mathbf{G}\) 中的）是下列映射的复合：映射 \(((x_i, y_i)) \to (x_i y_i^{-1})\)（从

\[
\prod_ {i \in I} G _ {i} \times G _ {i} \text {   into   } G,
\]

）以及典范映射 \(((x_{i}), (y_{i})) \to ((x_{i}, y_{i}))\)（从

\[
\mathbf {G} \times \mathbf {G} \quad \text { onto } \quad \prod_ {i \in \mathbf {I}} (\mathbf {G} _ {i} \times \mathbf {G} _ {i});
\]

）；而这两个映射都是连续的（第一章 § 4 第 1 目命题 1 的推论 1，以及命题 2）。

定义 2 把积集

\[
\mathbf {G} = \prod_ {i \in \mathbf {I}} \mathbf {G} _ {i}
\]

赋予诸 \(G_{i}\) 的群结构的积这个群结构与诸 \(G_{i}\) 的拓扑的积这个拓扑所得的拓扑群，称为诸拓扑群 \(G_{i}\) 的积。

若 \((\mathbf{J}_{\kappa})_{\kappa \in \mathbb{K}}\) 是 I 的一个划分，则 G 同构于诸拓扑群 \(\prod_{i \in J_{\kappa}} G_{i}\) 的积（积的结合性）。

若 \(\mathbf{H}_{i}\) 是 \(\mathbf{G}_{i}\) 的一个子群，则诸拓扑群 \(\mathbf{H}_{i}\) 的积同构于子群 \(\prod_{i} \mathbf{H}_{i}\)（\(\prod_{i} \mathbf{G}_{i}\) 的）。特别地，若 \(J\) 是 \(I\) 的任一子集，且 \(J' = C_{J}\)，则拓扑群 \(\prod_{i \in J} \mathbf{G}_{i}\) 同构于正规子群 \(\mathbf{G}_{J}' = (\prod_{i \in J} \mathbf{G}_{i}) \times (\prod_{i' \notin J} \{e_{i'}\})\)（\(\mathbf{G}\) 的）。由于每个开集的投影是开集，投影 \(\mathbf{pr}_{J}\)（\(\mathbf{G}\) 到 \(\prod_{i \in J} \mathbf{G}_{i}\) 上的）是严格态射，因而商群 \(\mathbf{G} / \mathbf{G}_{J}'\) 同构于 \(\mathbf{G}_{J'}'\)；\(\mathbf{G}\) 同构于积 \(\mathbf{G}_{J}' \times (\mathbf{G} / \mathbf{G}_{J}')\)。

命题 25 设 \((\mathbf{G}_{i})_{i\in I}\) 是一个拓扑群族，H 是 \(\mathbf{G} = \prod_{i\in I}\mathbf{G}_i\) 的由所有这样的 \(x = (x_{i})\) 组成的正规子群：诸 \(x_{i}\) 除有限个指标外都等于单位元 \(e_i\)（\(\mathbf{G}_i\) 的）。则子群 H 在 G 中稠密。

这是第一章 § 4 第 3 目命题 8 的一个特殊情形。

设 \((\mathbf{X}_i)_{i \in I}\) 是一个拓扑空间族，且对每个 \(i \in I\)，设 \(\mathbf{G}_i\) 是连续作用在 \(\mathbf{X}_i\) 上的一个拓扑群。显然，积群 \(\mathbf{G} = \prod_{i \in I} \mathbf{G}_i\) 这时依照法则连续作用在积空间 \(\mathbf{X} = \prod_{i \in I} \mathbf{X}_i\) 上

\[
\left(\left(s _ {i}\right), \left(x _ {i}\right)\right)\rightarrow \left(s _ {i}. x _ {i}\right)
\]

（第一章 § 4 第 1 目命题 1 推论 1，以及命题 2）。此外，X 的点 \(x = (x_{i})\) 在 G 下的轨道是诸 \(x_{i}\) 的轨道（关于诸群 \(G_{i}\) 的）的积。设 \(\varphi_{i}\) 是 \(X_{i}\) 到 \(X_{i}/G_{i}\) 上的典范映射，\(\varphi = (\varphi_{i})\) 是 X 到 \(\prod_{i\in I}(X_{i}/G_{i})\) 上的积映射；则前面的注记表明，典范地伴随于 \(\varphi\) 的双射把轨道空间 X/G 映为 \(\prod_{i\in I}(X_{i}/G_{i})\)。此外：

命题 26 X/G 到 \(\prod_{i\in I}(\mathbf{X}_i / \mathbf{G}_i)\) 上的双射（典范地伴随于 \((\varphi_i)\) 的）是一个同胚。

因为诸 \(\varphi_{i}\) 是满射且开的，可知 \(\varphi = (\varphi_{i})\) 是一个开映射（第一章 § 5 第 3 目命题 8 的推论）。

推论 设 \((\mathbf{G}_{i})_{i\in I}\) 是一个拓扑群族，且对每个 \(i\in I\)，设 \(\mathbf{H}_i\) 是 \(\mathbf{G}_i\) 的一个正规子群；以 \(\varphi_i\) 表示 \(\mathbf{G}_i\) 到 \(\mathbf{G}_i / \mathbf{H}_i\) 上的典范映射。设 \(\mathbf{G} = \prod_{i\in I}\mathbf{G}_i\)，\(\mathbf{H} = \prod_{i\in I}\mathbf{H}_i\)。则 \(\mathbf{G} / \mathbf{H}\) 到 \(\prod_{i\in I}(\mathbf{G}_i / \mathbf{H}_i)\) 上的、伴随于连续同态 \((x_i)\to (\varphi_i(x_i))\) 的双射同态是拓扑群的一个同构。

因为这个同态是群结构的一个同构。

注。若 G 是一个交换拓扑群，以加法记，则映射 \((x,y)\to x+y\)（\(G\times G\) 到 G 上的）是一个严格态射。因为由于 \((x+x')+(y+y')=(x+y)+(x'+y')\)，它是 \(G\times G\) 到 G 上的一个同态；它也是连续的，且邻域 \(V\times V\)（\(G\times G\) 中原点的）在这个映射下的像是 G 中原点的邻域 \(V+V\)。

